\section{Scope, limitations, and falsifiable predictions}
\label{sec:limits}

We state the boundaries of the present results explicitly; two of them are the outcome of
tests that could have gone either way, and both are reported honestly.

\paragraph{The carrier is a texture, not the ground state.} The helical carrier
of Section~\ref{subsec:carrier-construct} is linearly (elastically) stable---every link Hessian
bracket is positive-definite (minimum $1-r/\bar L=+0.20$), so the spatial fluctuation Hessian is
positive semidefinite and all fluctuation frequencies obey $\omega^2\ge 0$ \citep{repo}. However
it stores $+42.5\%$ more elastic energy per node than the straight vacuum: it is a stable
finite-energy \emph{texture}, not the elastic ground state. Since the straight ground state has
no color (trivial bundle, Section~\ref{subsec:stiffness}), the defensible claim is that an
$\su{3}$ sector \emph{can live on a stable substrate texture}, not that the ground-state vacuum
is a Yang--Mills field. Genuine (Floquet) selection of the vacuum carrier---and whether a
color-carrying background at or below the straight-vacuum energy exists---is open. (The
indefinite action is a saddle, so energy minimization does not by itself decide the vacuum.)

\paragraph{Emergent Lorentz invariance is approximate.} As shown in Section~\ref{sec:metric},
the lab dispersion is genuinely anisotropic for $\alpha_s>0$. We tested whether the physically
relevant gauge (Berry) modes are more isotropic than the bare phonons: the orientation-averaged
gauge kinetic tensor has anisotropy $\approx 0.81$, \emph{larger} than the bare acoustic
$\approx 0.25$ \citep{repo}. This candidate resolution is therefore falsified; the obstruction
persists in the gauge sector, and full cross-sector Lorentz universality remains the load-bearing
dual-observer conjecture, with only the matter-scale and observer-renormalization routes left
open. Residual birefringence ($c_L\neq c_T$, and direction-dependent transverse speeds) is a
concrete falsifiable signature.

\paragraph{This approximate Lorentz invariance is a feature: the Weinberg--Witten escape.}
The Weinberg--Witten theorem \citep{WeinbergWitten1980} forbids emergent massless charged spin-1
(and spin-2) particles in a theory with an \emph{exactly} Lorentz-covariant current---precisely
the case of an emergent photon and gluon. The standard escape, shared with the emergent-gravity
and emergent-gauge literature \citep{Volovik2003,Wen2004}, is that Lorentz invariance is only
approximate, so the exact-covariance premise fails. The residual anisotropy established above is
therefore not merely a limitation but the property that permits the emergent $U(1)\times SU(3)$
sector to coexist with the no-go theorem: the gauge bosons are collective substrate modes, and
exact Lorentz covariance holds only in the long-wavelength limit, with corrections at
$O((\veck a)^2)$ and $O(\alpha_s)$.

\paragraph{Matter, confinement, and dynamics are deferred.} No self-confined solitons are
constructed here; hence no masses, no color sources $J^a_\mu$, no confinement, and no contact
with measured couplings or mass ratios. Confinement and the mass gap are nonperturbative and are
not addressed. The absolute normalizations of $1/e^2$ and $1/g^2$, and the exact
$\beta$-function, involve the coarse-graining scheme and are left open; only positivity,
finiteness, and scale dependence are established. The Minkowski index contractions used in the
effective actions inherit the emergent metric, whose universality is the conjecture above.

\paragraph{Foundational open problems.} The well-posedness of the nonlinear two-time (block)
boundary-value problem is open (it is resolved only in the linear regime, via a chiral two-slice
condition). The identification of the trace $U(1)$ with the physical electromagnetic field is a
candidate pending the matter sector. Quantum-mechanical probability (the Born rule) is not derived.
